\Askhsh{34151}{4}
\begin{ylh}{1}
    \item 2.2 Παραγωγίσιμες συναρτήσεις - Παράγωγος συνάρτηση
    \item 2.3 Κανόνες παραγώγισης
    \item 2.6 Συνέπειες του Θεωρήματος Μέσης Τιμής
    \item 2.7 Τοπικά ακρότατα συνάρτησης
    \item 3.5 Η συνάρτηση [ορισμένο ολοκλήρωμα της f από α έως χ]
    \item 3.7 Εμβαδόν επιπέδου χωρίου.
\end{ylh}
\begin{wrapfigure}[20]{r}{9cm}
    \begin{tikzpicture}
        \begin{axis}[
            axis lines = middle,
            xlabel = {$x$},
            ylabel = {$y$},
            xmin = -1.5, xmax = 2.5,
            ymin = -1.75, ymax = 3.5,
            xtick = {-1.5, -1, -0.5, 0, 0.5, 1, 1.5, 2, 2.5},
            ytick = {-1.5, -1, -0.5, 0, 0.5, 1, 1.5, 2, 2.5, 3},
            grid = both,
            grid style = {gray!30},
            samples = 200,
            width = 8.5cm, height = 6cm,
            tick label style = {font=\scriptsize},
            every axis x label/.style = {at={(current axis.right of origin)}, anchor=north west},
            every axis y label/.style = {at={(current axis.above origin)}, anchor=south},
        ]
            % C2: f(x) = 4x^3 - 6x^2 + 2x (blue -> maincolor)
            \addplot[very thick, maincolor, domain=-0.2:1.2] {4*x^3 - 6*x^2 + 2*x};
            \node[maincolor, right] at (axis cs:1.3,-0.1) {$C_2$};

            % C1: f'(x) = 12x^2 - 12x + 2 (red -> secondarycolor)
            \addplot[very thick, secondarycolor, domain=-0.2:1.2] {12*x^2 - 12*x + 2};
            \node[secondarycolor, left] at (axis cs:0.1,2.8) {$C_1$};

            % C3: F(x) = x^4 - 2x^3 + x^2 + 1 (green -> thirdcolor)
            \addplot[very thick, thirdcolor, domain=-0.2:1.2] {x^4 - 2*x^3 + x^2 + 1};
            \node[thirdcolor, above left] at (axis cs:-0.3,3) {$C_3$};
        \end{axis}
    \end{tikzpicture}
\end{wrapfigure}
Στο παρακάτω σχήμα δίνονται οι γραφικές παραστάσεις $C_1, C_2, C_3$ τριών συναρτήσεων $f', f$ και $F$, όπου $F$ μία αρχική της $f$ στο $\mathbb{R}$. Δίνεται επίσης ότι η $C_3$ τέμνει τον άξονα $y'y$ στο σημείο με τεταγμένη $1$ ενώ η $C_2$ διέρχεται από την αρχή των αξόνων και τέμνει τον άξονα $x'x$ σε δύο ακόμη σημεία με τετμημένες $\frac{1}{2}, 1$. Με δεδομένο ότι ο τύπος της $f$ είναι $f(x) = 4x^3 - 6x^2 + 2x$ και η γραφική της παράσταση είναι η $C_2$.
\begin{alist}
    \item να μελετήσετε, με τη βοήθεια του σχήματος ή με οποιονδήποτε άλλο τρόπο, τη συνάρτηση $F$ ως προς την μονοτονία και τα ακρότατα. \monades{7}
    \item να δικαιολογήσετε γιατί η γραφική παράσταση $C_3$ αντιστοιχεί στην συνάρτηση $F$. \monades{6}
    \item να βρείτε τον τύπο των συναρτήσεων $f'$ και $F$. \monades{6}
    \item να βρείτε το εμβαδόν του χωρίου που περικλείεται μεταξύ του άξονα $x'x$ και της γραφικής παράστασης της συνάρτησης $f'$. \monades{6}
\end{alist}